% Physical pages 3 and 4 of the concise study: the thematic orientation, in
% the order of the Mass, with the commentators' answers to the questions
% each text raises. Every claim and quotation comes from the reviewed
% expansive study. The developed comparison of the three interpretations,
% organized by the questions the texts raise, begins on the fifth page.

\runninghead{Themes and Movement}
\section{The Propers: Themes and Movement}
\zlabel{triptych:concise:themes:start}

The Mass of All Saints opens with words of the Church, not of Scripture:
``Let us all rejoice in the Lord, celebrating a feast in honor of all the
saints, in whose solemnity the angels rejoice and join in praising the Son
of God.'' From that joy the Mass moves to John's vision of the servants of
God, counted and then uncountable, to the psalm's promise that those who
fear the Lord lack nothing, to the Beatitudes, the way the saints walked,
and to the just who only seemed to die. Its prayers ask throughout for the
intercession of the saints who have arrived, on behalf of the faithful
still on the way.

\subsection{Introit and Collect: one celebration of all the saints}

The Introit antiphon was not written for this feast. Schuster notes that
\latin{Gaudeamus \dots\ sub honore} is an ancient Introit ``originally
assigned to the feast of St Agatha on February 5,'' and Durandus already
gives it to her, explaining it by the joy with which she went to prison and
to her passion (\work{The Sacramentary} V, p.~209; \work{Rationale} VII.47).
On this feast the antiphon names \latin{Sanctorum omnium}, all the saints,
and keeps the angels who rejoice at the solemnity. The verse opens
Psalm~32, and St Augustine hears in it a caution about where the joy rests:
``rejoice, O ye righteous, not in yourselves, for that is not safe; but in
the Lord'' (\work{En.\ in Ps.}\ 32 [NPNF ``Psalm XXXIII''], 1).

The Collect says that God has granted us to venerate the merits of all his
saints \latin{sub una celebritate}, in one celebration, and asks for the
abundance of his mercy \latin{multiplicatis intercessoribus}, the
intercessors being multiplied. Durandus, giving the causes for which the
feast was instituted, names both thoughts: many saints have no feast of
their own, and petitions are more easily granted when all the saints
intercede together, \latin{quia impossibile est multitudinem non exaudiri},
because it is impossible that a multitude should not be heard
(\work{Rationale} VII.34). After it the priest says the commemorated
Sunday's Collect, that God loose the sins of his peoples from the bonds
their frailty contracted.

\subsection{Lesson: a number, and a multitude without number}

The Lesson (Apoc 7:2--12) holds the feast's first question. An angel rising
from the east forbids the four angels who have power to hurt the earth to
hurt nothing until the servants of God are signed on their foreheads. John
hears the number of the sealed, a hundred and forty-four thousand, given
again tribe by tribe, twelve thousand from each of twelve --- and the list
is not the list of Jacob's sons in the order of their birth: Juda stands
first, Dan is missing, Manasses is counted beside his father Joseph. Then
John sees ``a great multitude, which no man could number, of all nations and
tribes and peoples and tongues,'' in white robes and with palms, before the
throne and the Lamb.

St Bede reads the angel as the Lord incarnate, ``the day-spring from on
high,'' bearing the ensign of the cross with which to seal his own; the
definite number, he says, signifies ``the innumerable multitude of the whole
Church, which is descended from the patriarchs either by the lineage of
nature, or the imitation of faith,'' for ``if ye be Christ's, then are ye
Abraham's seed'' (Gal 3:29) (\work{Explanatio Apocalypsis} I, on 7:2--4).
The robes are baptism, and the palms ``the triumph of the Cross,'' for ``in
Christ they have overcome the world'' (on 7:9). St Augustine, who expounds
the chapter nowhere, says in passing that the number of the Apocalypse
signifies the whole body of the saints (\work{De doctrina christiana}
III.35.51).

\subsection{Gradual, Alleluia and Gospel: the way and its promise}

The Gradual sings Psalm~33: ``Fear the Lord, all ye his saints: for there
is no want to them that fear him,'' and then that those who seek the Lord
shall lack no good; the first half of the verse, ``The rich have wanted,
and have suffered hunger,'' is not sung. St Basil preaches that the fear
which orders a life lacks nothing for any virtue, and that the wholly
perfect good is God himself, \latin{Perfectum omnino bonum Deus ipse est}
(\work{Hom.\ in Ps.}\ 33, PG 29, cols.~365--368). St Robert Bellarmine reads
the promise in a higher sense, in which the friend of God who owns
everything feels no want, and in a humbler one, in which God does not
desert his servants in their necessity (\work{Explanatio}, Ps~33, his
v.~9). The Alleluia is the Lord's invitation, ``Come to me all you that
labor and are burdened, and I will refresh you'' (Mt 11:28), which
Chrysostom notes reaches ``not this or that person, but all that are in
anxiety, in sorrows, in sins'' (\work{Hom.\ in Matt.}\ 38.3).

The Gospel is the opening of the Sermon on the Mount (Mt 5:1--12a), which
St Augustine calls ``a perfect standard of the Christian life'' (\work{De
serm.\ Dom.}\ I.1). He reads the eight beatitudes as seven steps, from the
fear of the Lord to wisdom, with an eighth that ``returns to the
starting-point, because it shows and commends what is complete and
perfect''; the seven answer to the sevenfold operation of the Holy Ghost in
Isaias 11, and the rewards are one reward, the kingdom of heaven,
``variously named according to these stages'' (I.3--4). St John Chrysostom
reads ``a sort of golden chain'': he that is humble will mourn, he that
mourns will be meek and righteous and merciful, ``and he that hath attained
unto all these, will be moreover arrayed against dangers'' (\work{Hom.\ in
Matt.}\ 15.9). St Thomas Aquinas answers the question of tense: the rewards
are had perfectly only in the homeland, and in their beginning on the way,
and the poor in spirit are blessed \latin{non spe tantum, sed etiam re},
not only in hope but in fact (\work{Super Matt.}\ cap.~5, lect.~1). St
Hilary of Poitiers reads the promises of the life to come: the clean of
heart will see God when they are perfected in spirit and changed by
immortality (\work{In Matt.}\ IV).

\subsection{Offertory to Postcommunion: in the hand of God}

The Offertory adapts Wisdom's answer to the wicked, who had resolved to
condemn the just man to a shameful death (Wis 2:20; 3:1--3): ``The souls of
the just are in the hand of God, and the torment of malice shall not touch
them: in the sight of the unwise they seemed to die, but they are in
peace.'' St Gregory the Great applies the verse to the just man whose body
is left lifeless ``amidst the hands of the persecutors'' (\work{Moralia}
XII.v.6), and Schuster states the contrast in terms of hands: ``The
persecutor believes that the lives of the martyrs are in his hands, but
they are in the hands of God'' (p.~211).

The Secret asks that the gifts of our devotion be pleasing ``for the honor
of all the just'' and, by God's mercy, salutary to those who offer them;
the commemorated Sunday's Secret, said after it, asks God to bring to
completion what he has given to the undeserving. The Communion sings again
three of the beatitudes the Gospel has just given --- the clean of heart,
the peacemakers and the persecuted --- whose promises name the way's end:
to see God, to be called the children of God, to possess the kingdom of
heaven. The Postcommunion asks that the \latin{fideles populi}, the
faithful peoples, ever rejoice in the veneration of all the saints and be
defended by their unceasing prayer; Schuster hears in that name the Church
on earth, ``journeying towards eternity with the eyes, and in the light of
faith'' (p.~211). The Sunday's Postcommunion, the last proper prayer of the
Mass, asks that those whom God gives to rejoice in sharing divine things be
not left subject to human dangers.
\zlabel{triptych:concise:themes:end}
